\section{Spektrogram i ResNet2D}

\subsection{Spektrogram}
Domyślnie sygnał EKG jest jednowymiarowy. Lekarze jednak korzystają czasem ze spektrogramu, 
który pokazuje, jak zmienia się zawartość sygnału w czasie. Dzięki temu z jednowymiarowego przebiegu $(12, L)$
otrzymujemy dwuwymiarową mapę $(12, F, T)$, gdzie $F$ jest liczbą pasm częstotliwości, a $T$
to liczba okien czasowych.

Spektrogram obliczamy przy użyciu STFT,
korzystając z gotowej implementacji dostępnej w bibliotece $scipy$.
Przesuwa ona krótkie okno czasowe po sygnale i dla każdej pozycji sprawdza,
jakie częstotliwości w nim dominują. W naszym przypadku okno ma długość 128 próbek
i przesuwa się co 32 próbki. Wyniki w powstałej macierzy logarytmujemy, żeby pozornie
małe różnice między niskimi wynikami nie zostały "wyciszone" względem wyników wysokich.

\subsection{ResNet2D}
Mając spektrogram jako wejście w postaci obrazu $(12, F, T)$, możemy zastosować
standardową sieć konwolucyjną operującą w dwóch wymiarach. Wykorzystujemy tę
samą architekturę ResNet18 opisaną w poprzedniej sekcji, z tą różnicą, że
wszystkie warstwy konwolucyjne są dwuwymiarowe.

To podejście jest o tyle unikalne, że sieć może jednocześnie uczyć się
wzorców zarówno w osi czasu, jak i w osi częstotliwości, co jest trudniejsze
do osiągnięcia przy pozostałych modelach.
